\documentclass[11pt,a4paper]{article}
\usepackage[margin=1in]{geometry}
\usepackage{graphicx}
\usepackage{booktabs}
\usepackage{amsmath}
\usepackage{siunitx}
\usepackage{float}
\usepackage{hyperref}
\usepackage{caption}
\usepackage{subcaption}

\title{Reproduction Report: Optimization-Driven Quantum Circuit Reduction}
\author{Arnav Kapoor 23060}
\date{\today}

\begin{document}
\maketitle

\begin{abstract}
In this report, I document not only what I reproduced from \emph{Optimization driven quantum circuit reduction} but also how I made implementation decisions under real constraints (tooling, data format, and verification fidelity). I translated the MATLAB demo flow into a Python workflow, built reproducible benchmarks, generated figures, and added baseline/future-work artifacts. On the 5-qubit demo circuit (300 initial gates), my best strict run reached 228 gates (depth 4, 1500 iterations), while a loose-tolerance run reached 226 gates. The report is written as a technical workflow narrative to demonstrate my research process and readiness for lab work.
\end{abstract}

\section{Motivation and objective}
My goal was to show that I can take a research paper, operationalize it in code, and verify claims with reproducible evidence. I focused on the paper's core idea: reducing transpiled circuit length using local replacement while preserving unitary action (up to global phase). I treated this as a mini research project rather than a script-only exercise: I built infrastructure, handled blockers, tested assumptions, and documented outcomes.

\section{My workflow (what I did and why)}
\subsection{Stepwise execution log}
I followed the workflow in Table~\ref{tab:workflow}. Each step includes both action and rationale.

\begin{table}[H]
\centering
\caption{My reproduction workflow and rationale.}
\label{tab:workflow}
\begin{tabular}{p{0.13\linewidth}p{0.39\linewidth}p{0.41\linewidth}}
\toprule
Phase & What I did & Why I did it \\
\midrule
1 & Parsed the paper excerpt and extracted algorithmic requirements (V1/V2/V3, tolerances, gate sets). & I needed a testable specification before writing any code. \\
2 & Set up a clean Python environment and project scaffold. & Reproducibility requires deterministic dependencies and scripted execution. \\
3 & Integrated the provided archive \texttt{QCOptimDemo.zip} and inspected MATLAB scripts. & The archive is the closest available ground-truth implementation reference. \\
4 & Ported the MATLAB demo logic to Python (QASM parsing, local remap, compute-graph lookup, replacement loop). & I wanted an executable, inspectable baseline in a modern workflow. \\
5 & Added equivalence checks and ran strict vs loose tolerance validation. & Numeric tolerance strongly affects whether reductions are accepted. \\
6 & Built benchmark sweeps, generated figures/tables, and produced a report + PDF. & I needed evidence quality comparable to a lab handoff. \\
\bottomrule
\end{tabular}
\end{table}

\subsection{Artifacts and environment}
I prepared the following assets:
\begin{itemize}
    \item Reference code archive extracted to \texttt{paper\_code/QCOptimDemo} (MATLAB scripts and data files).
    \item Python virtual environment at \texttt{.venv} with dependencies in \texttt{requirements.txt}.
    \item Python implementation under \texttt{src/qcr\_repro} and experiment scripts under \texttt{scripts}.
\end{itemize}

\subsection{Technical implementation decisions}
The implemented reduction pipeline follows the paper's local replacement concept:
\begin{enumerate}
    \item Parse QASM circuit into gate tokens.
    \item Sample local contiguous blocks.
    \item Remap active wires to local coordinates (MATLAB \texttt{ExtractWireMap}/\texttt{WireMapF} analogue).
    \item Query precomputed local compute-graph lookup.
    \item Accept replacement only when (i) shorter and (ii) unitary-consistent.
\end{enumerate}
For practical parity with the demo code, I used local operator discretization
$\mathrm{RX},\mathrm{RY},\mathrm{RZ}\in\{\pm\pi/2\}$ and $\mathrm{RXX}(\pi/2)$.

\subsection{Key debugging and engineering choices}
Two implementation choices had major impact:
\begin{itemize}
    \item \textbf{Compute cost control:} depth-4 lookup can become expensive. I reduced operator pool entropy to match demo discretization and introduced commutation memoization to avoid repeated expensive checks.
    \item \textbf{Validation policy:} I evaluated both strict ($10^{-5}$) and loose ($10^{-3}$) tolerances, because rounded QASM constants in demo files can fail strict numerical equivalence despite being practically close.
\end{itemize}

\section{Results}
\subsection{Benchmark protocol}
I benchmarked over depths $\{3,4\}$, iteration budgets $\{500,1000,1500\}$, and seeds $\{1,5,10\}$ (18 runs per tolerance setting).

\subsection{Best configurations}
\begin{itemize}
    \item Best strict-valid run ($10^{-5}$): depth 4, iterations 1500, seed 10, end length 228, runtime 29.8365 s.
    \item Best loose-valid run ($10^{-3}$): depth 4, iterations 1500, seed 5, end length 226, runtime 2.4003 s.
\end{itemize}

Interpretation: the loose metric improves acceptance and runtime, while strict verification is more conservative and exposes fragile reductions.

\subsection{Grouped benchmark table}
\begin{table}[H]
\centering
\caption{Combined strict vs loose tolerance results (3 seeds per configuration).}
\label{tab:combined}
\begin{tabular}{rrrrrrrrr}
\toprule
depth & iters & runs & best$_{strict}$ & mean$_{strict}$ & eq$_{strict}$ & best$_{loose}$ & mean$_{loose}$ & eq$_{loose}$ \\
\midrule
3 & 500  & 3 & 266 & 274.00 & 1.000 & 266 & 274.00 & 1.000 \\
3 & 1000 & 3 & 246 & 250.00 & 0.667 & 246 & 250.00 & 1.000 \\
3 & 1500 & 3 & 240 & 242.67 & 0.667 & 240 & 242.67 & 1.000 \\
4 & 500  & 3 & 262 & 265.33 & 1.000 & 262 & 265.33 & 1.000 \\
4 & 1000 & 3 & 238 & 243.33 & 1.000 & 238 & 243.33 & 1.000 \\
4 & 1500 & 3 & 226 & 231.33 & 0.667 & 226 & 231.33 & 1.000 \\
\bottomrule
\end{tabular}
\end{table}

\section{Figure-based analysis}
I generated all implementation figures into \texttt{implementation/figures}. They summarize my empirical workflow:
\begin{itemize}
    \item Figure~\ref{fig:motivation}: reduction endpoints (original, MATLAB short, Python strict/loose).
    \item Figure~\ref{fig:graphgrowth}: compute-graph growth behavior.
    \item Figure~\ref{fig:pipeline}: algorithmic pipeline I implemented.
    \item Figure~\ref{fig:curve}: convergence trend with increasing iteration budget.
    \item Figure~\ref{fig:runtime}: runtime--quality tradeoff.
    \item Figure~\ref{fig:boxplot}: variance across seeds under strict checks.
\end{itemize}

\begin{figure}[H]
    \centering
    \includegraphics[width=0.85\linewidth]{../figures/figure1_motivation.png}
    \caption{Implementation Figure 1: circuit-length comparison between original and reduced variants.}
    \label{fig:motivation}
\end{figure}

\begin{figure}[H]
    \centering
    \includegraphics[width=0.95\linewidth]{../figures/figure2_compute_graph_growth.png}
    \caption{Implementation Figure 2: compute-graph growth with depth for representative operator pools.}
    \label{fig:graphgrowth}
\end{figure}

\begin{figure}[H]
    \centering
    \includegraphics[width=0.95\linewidth]{../figures/figure3_pipeline.png}
    \caption{Implementation Figure 3: reduction pipeline (sampling, lookup, replacement).}
    \label{fig:pipeline}
\end{figure}

\begin{figure}[H]
    \centering
    \includegraphics[width=0.85\linewidth]{../figures/figure4_reduction_curve.png}
    \caption{Implementation Figure 4: reduction trend versus iteration budget.}
    \label{fig:curve}
\end{figure}

\begin{figure}[H]
    \centering
    \includegraphics[width=0.85\linewidth]{../figures/figure5_runtime_vs_length.png}
    \caption{Implementation Figure 5: runtime versus achieved end length across benchmark runs.}
    \label{fig:runtime}
\end{figure}

\begin{figure}[H]
    \centering
    \includegraphics[width=0.95\linewidth]{../figures/figure6_boxplot.png}
    \caption{Implementation Figure 6: distribution of end lengths across seeds (strict setting).}
    \label{fig:boxplot}
\end{figure}

\section{Discussion}
The core behavior reported in the paper is reproduced: local stochastic search plus lookup replacement reduces circuit length substantially. More importantly, this project highlighted a realistic research lesson: \emph{validation criteria drive conclusions}. Under strict checks, some reductions are rejected; under loose checks, acceptance improves and trends become cleaner. For my workflow, this meant keeping both views side-by-side and explicitly reporting the tradeoff instead of hiding it.

\section{What this demonstrates about my research workflow}
I want to explicitly state what I think this work shows:
\begin{itemize}
    \item I can convert a paper plus partial artifacts into a working, benchmarked implementation.
    \item I can diagnose bottlenecks (runtime/memoization/tolerance sensitivity) and patch them systematically.
    \item I can produce reproducible outputs (scripts, figures, tables, and a compiled report) suitable for lab collaboration.
\end{itemize}

\section{Future-work implementation in folder}
I implemented the requested future-work direction in \texttt{implementation/future\_work}:
\begin{itemize}
    \item \texttt{baseline\_parity\_results.csv}: baseline parity runs for local method and Qiskit optimization levels.
    \item \texttt{hardware\_metrics.json}: hardware-level metric probe output.
    \item \texttt{future\_work\_status.md}: concise status summary.
\end{itemize}

The recorded baseline parity results are:
\begin{itemize}
    \item Local method (strict): 228 gates in 2.4156 s, equivalence check passed.
    \item Qiskit transpile baseline (same gate basis):
    \begin{itemize}
        \item optimization level 0: 300 gates,
        \item optimization level 1: 229 gates,
        \item optimization level 2: 219 gates,
        \item optimization level 3: 219 gates.
    \end{itemize}
    \item BQSKit: package availability verified; full pass-level parity script remains as next incremental step.
\end{itemize}

For hardware-level execution metrics, no authenticated hardware provider was available in this environment, and this status is explicitly logged in \texttt{hardware\_metrics.json}.

\section{Remaining limitations}
Two constraints remain for full paper-level parity:
\begin{itemize}
    \item The supplied archive is MATLAB-focused demo code; my main pipeline is Python-first for reproducibility.
    \item Full replication of all paper claims (especially hardware experiments) still depends on provider access and complete backend details.
\end{itemize}

\section*{Reproducibility note}
All generated outputs used in this report are present in the workspace under \texttt{results}, \texttt{results\_tol1e3}, \texttt{implementation/figures}, and \texttt{implementation/future\_work}.

\end{document}
